% Fig.Intro: LLM-to-MCU pipeline overview (IMPROVED VERSION)
% 改进点:
% 1. 字体加大: \scriptsize → \footnotesize (节点), \tiny → \scriptsize (标注)
% 2. 节点尺寸: 0.85cm×1.8cm → 1.2cm×2.6cm
% 3. 箭头加粗: thick → very thick
% 4. 层背景加深: black!3 → blue!5, 边框 black!15 → black!35
% 5. 间距加大: 3.2cm → 4.0cm (消除挤压)
% 6. 反馈箭头突出: 加粗至1.5pt, 改为红色, 更大弧线
% 7. 关键节点高亮: LLM/LVM用深色填充
% 8. 对齐优化: 严格grid布局

\begin{tikzpicture}[scale=1.0, transform shape,
  >=Stealth, font=\sffamily\small,
  % 节点样式 - 加大尺寸,加大字体
  box/.style={draw=black!70, fill=white, minimum width=2.6cm,
              minimum height=1.2cm, align=center, inner sep=4pt,
              font=\sffamily\footnotesize\bfseries, line width=0.6pt},
  % 高亮节点 - 关键组件(LLM/LVM)
  boxhl/.style={draw=blue!80, fill=blue!25, minimum width=2.6cm,
              minimum height=1.2cm, align=center, inner sep=4pt,
              font=\sffamily\footnotesize\bfseries, line width=0.8pt},
  % 警告节点 - 安全关键(Bytecode)
  boxwarn/.style={draw=orange!80, fill=orange!15, minimum width=2.6cm,
              minimum height=1.2cm, align=center, inner sep=4pt,
              font=\sffamily\footnotesize\bfseries, line width=0.7pt},
  % 箭头 - 加粗
  arr/.style={-Stealth, very thick, draw=black!70},
  % 反馈箭头 - 突出显示
  feedback/.style={-Stealth, line width=1.5pt, draw=red!75!black, dashed},
  % 标注文字 - 加大字号
  note/.style={font=\sffamily\scriptsize, black!60, align=center},
]

% ============================================================
% Layer 1: Generation-friendly (顶层)
% ============================================================
% 背景框 - 加深颜色,加粗边框
\fill[blue!5, draw=black!35, line width=0.5pt, rounded corners=2pt]
  (-0.8, 2.0) rectangle (16.8, 4.8);
\node[font=\sffamily\large\bfseries, black!70, anchor=north west]
  at (-0.7, 4.7) {Generation-friendly Layer};

% 节点 - 横向间距4.0cm,纵向居中y=3.4
\node[box]    (nl)   at (1.2, 3.4) {Natural\\Language};
\node[boxhl]  (llm)  at (5.2, 3.4) {LLM};
\node[box]    (lir)  at (9.2, 3.4) {LIR\\(DSL)};
\node[box]    (cmp)  at (13.2, 3.4) {Compiler};

% 标注 - 加大字号,下移避免与节点贴太近
\node[note] at (1.2, 2.55) {user\\instruction};
\node[note] at (5.2, 2.55) {text\\generation};
\node[note] at (9.2, 2.55) {EBNF grammar\\LLM-stable};
\node[note] at (13.2, 2.55) {deterministic\\lowering};

% 箭头 - 加粗
\draw[arr] (nl.east) -- (llm.west);
\draw[arr] (llm.east) -- (lir.west);
\draw[arr] (lir.east) -- (cmp.west);

% ============================================================
% Layer 2: Execution-friendly (底层)
% ============================================================
% 背景框
\fill[orange!5, draw=black!35, line width=0.5pt, rounded corners=2pt]
  (-0.8, -0.7) rectangle (16.8, 1.8);
\node[font=\sffamily\large\bfseries, black!70, anchor=north west]
  at (-0.7, 1.7) {Execution-friendly Layer};

% 节点 - 对齐上层
\node[boxwarn] (bc)  at (5.2, 0.6) {Compact\\Bytecode};
\node[boxhl]   (lvm) at (9.2, 0.6) {LVM\\($\sim$2\,KB)};
\node[box]     (dev) at (13.2, 0.6) {MCU\\Device};

% 标注
\node[note] at (5.2, -0.25) {10.9\,B avg\\13.8$\times$ vs JSON};
\node[note] at (9.2, -0.25) {bounded exec\\UABR\,=\,1.000};
\node[note] at (13.2, -0.25) {relay/sensor\\ESP8266};

% 箭头
\draw[arr] (bc.east) -- (lvm.west);
\draw[arr] (lvm.east) -- (dev.west);

% ============================================================
% 垂直连接: Compiler → Bytecode
% ============================================================
\draw[arr] (cmp.south) -- ++(0, -0.6) -| (bc.north)
  node[pos=0.4, right, font=\sffamily\scriptsize, black!65] {varint encode};

% ============================================================
% 反馈路径: LVM → LLM (repair loop)
% ============================================================
% 加粗红色虚线,更大弧度
\draw[feedback]
  (lvm.south) ++(0.8, 0) -- ++(0, -1.8)
  -| ([xshift=-0.8cm]llm.south)
  -- (llm.south)
  node[pos=0.08, below, font=\sffamily\small, red!75!black, align=center]
    {fault\\feedback};

% ============================================================
% 右侧约束标签 (设计目标)
% ============================================================
\node[font=\sffamily\small, black!50, anchor=west, text width=2.8cm, align=left]
  at (15.2, 4.1) {\textbullet~token-efficient\\[1pt]\textbullet~compact};
\node[font=\sffamily\small, black!50, anchor=west, text width=2.8cm, align=left]
  at (15.2, 0.6) {\textbullet~bounded-exec\\[1pt]\textbullet~auditable};

\end{tikzpicture}
